\subsection{Rough noise beyond the present covariance class}
\label{S:future-rough-noise-extension}

We ask whether we can construct the Stratonovich wave equation with initial
data $(1,0)$, and determine its moment growth, beyond
Hypothesis~\ref{H:homog}.
We use $d = 1$ to illustrate these questions and consider Gaussian noises with
product spectral measure $\mu_0\otimes\mu$, where
\[
  \mu_0(\ud\tau) = c_0|\tau|^{\alpha_0-1}\ud\tau,
  \qquad
  \mu(\ud\xi) = c_1|\xi|^{\alpha-1}\ud\xi,
  \qquad c_0, c_1 > 0,
\]
with $\alpha_0, \alpha > 0$ and $\alpha_0 + \alpha < 2$. The corresponding Hurst
parameters are $H_0 = 1 - \alpha_0/2$ and $H = 1 - \alpha/2$. Thus, the
proposed condition is $H_0 + H > 1$. We propose this region for investigation,
without asserting a solvability criterion. The
physical covariance in the rough regions is understood as a distribution and
need not be a nonnegative Radon measure.

After matching the covariance normalization, the current theorem covers
$0 < \alpha_0 < 1$, $0 < \alpha < 1$, and the spatial white noise segment
$\alpha = 1$, $0 < \alpha_0 < 1$.
Figure~\ref{Fig:rough-noise-region} separates this range from the two proposed
regions: rough space with $\alpha > 1$, and rough time with $\alpha_0 > 1$. The
temporal white noise segment $\alpha_0 = 1$, $0 < \alpha < 1$ is outside the
current theorem. Its integration convention and its relation to the present
variational formula require separate treatment. On either white noise line
we mean a Dirac covariance with
its specified normalization, not the literal function $|r|^{-1}$. The point
$(\alpha,\alpha_0) = (1,1)$ lies on the excluded boundary $\alpha_0 + \alpha =
2$; this exclusion is not a nonexistence statement.

\usetikzlibrary{patterns}
% Requires TikZ libraries arrows.meta and patterns.
\begin{figure}[htbp]
  \centering
  % Palette matches the book's hamblue/hamgold without changing shared styles.
\begingroup%
\definecolor{rnBlue}{RGB}{36,75,105}%
\definecolor{rnGold}{RGB}{156,101,20}%
\definecolor{rnGray}{RGB}{105,110,115}%
\begin{tikzpicture}[
  x=3.8cm,y=3.8cm,
  font=\small,
  region/.style={align=center,fill=white,fill opacity=.94,
  text opacity=1,inner sep=2pt,rounded corners=1pt},
  legend/.style={anchor=west,align=left,text width=4.4cm},
  axis/.style={-{Latex[length=2.1mm]},line width=.7pt},
  threshold/.style={draw=rnGray!45,densely dotted,line width=.8pt}
  ]
  % Three disjoint interiors within the proposed open triangle.
  \fill[rnBlue!16] (0,0) rectangle (1,1);
  \fill[rnGold!7] (1,0)--(2,0)--(1,1)--cycle;
  \fill[pattern=north east lines,pattern color=rnGold!38]
  (1,0)--(2,0)--(1,1)--cycle;
  \fill[rnGold!7] (0,1)--(1,1)--(0,2)--cycle;
  \fill[pattern=north west lines,pattern color=rnGold!38]
  (0,1)--(1,1)--(0,2)--cycle;

  % White-noise thresholds continue across the plot; exterior is subdued.
  \draw[threshold] (1,1)--(1,2.2);
  \draw[threshold] (1,1)--(2.2,1);
  \draw[rnBlue,line width=1.7pt] (1,0)--(1,1);
  \draw[rnGold,line width=1.35pt,dash pattern=on 5pt off 3pt]
  (0,1)--(1,1);
  \draw[rnGray,line width=1pt,dash pattern=on 4pt off 3pt]
  (0,2)--(2,0)
  node[pos=.25,above=6pt,sloped,text=rnGray]
  {$\alpha + \alpha_0 = 2$};

  \draw[axis] (0,0)--(2.2,0) node[right=5pt] {$\alpha$};
  \draw[axis] (0,0)--(0,2.2) node[above=5pt] {$\alpha_0$};
  \node[below left=3pt] at (0,0) {$0$};
  \foreach \v in {1,2} {
  \draw (\v,0)--++(0,-.025) node[below=4pt] {$\v$};
  \draw (0,\v)--++(-.025,0) node[left=4pt] {$\v$};
  }

  \node[region,text=rnBlue,text width=2.75cm] at (.49,.46)
  {\textbf{Current theorem}\\regular region\\
  $0 < \alpha < 1$, $0 < \alpha_0 < 1$};
  \node[region,text=rnGold,text width=2.4cm] at (1.35,.17)
  {\mbox{Rough space}\\\mbox{open question}};
  \node[region,text=rnGold,text width=2.4cm] at (.35,1.16)
  {\mbox{Rough time}\\\mbox{open question}};
  \node[text=rnGray,align=center,text width=3.0cm] at (1.67,1.79)
  {Rough in both\\time and space\\outside the proposed region};
  \node[text=rnGray,anchor=west,align=left,text width=3.0cm]
  at (1.13,1.20) {Both white $(1,1)$\\excluded};
  \draw[rnGray,line width=.6pt] (1.02,1.02)--(1.10,1.12);

  % Open markers: neither axes nor the diagonal belong to the range.
  \foreach \p in {(0,0),(0,1),(1,0),(0,2),(2,0)} {
  \filldraw[fill=white,draw=rnGray,line width=.7pt] \p circle[radius=2.2pt];
  }
  \filldraw[fill=white,draw=rnGray,line width=1.2pt]
  (1,1) circle[radius=3.4pt];

  % Legend uses fill, hatching and line style as well as color.
  \node[legend,font=\normalsize\bfseries,text=rnBlue] at (2.39,2.02)
  {Current theorem};
  \filldraw[fill=rnBlue!16,draw=rnBlue!50]
  (2.39,1.745) rectangle (2.57,1.855);
  \node[legend] at (2.66,1.80) {Regular region};
  \draw[rnBlue,line width=1.7pt] (2.39,1.52)--(2.57,1.52);
  \node[legend] at (2.66,1.52)
  {White noise in space\\$\alpha = 1$, $0 < \alpha_0 < 1$};

  \node[legend,font=\normalsize\bfseries,text=rnGold] at (2.39,1.20)
  {Extension questions};
  \filldraw[fill=rnGold!7,draw=rnGold!60]
  (2.39,.865) rectangle (2.57,.975);
  \fill[pattern=north east lines,pattern color=rnGold!55]
  (2.39,.865) rectangle (2.57,.975);
  \node[legend] at (2.66,.92) {Rough space: $\alpha > 1$\\rough time: $\alpha_0 > 1$};
  \draw[rnGold,line width=1.35pt,dash pattern=on 5pt off 3pt]
  (2.39,.58)--(2.57,.58);
  \node[legend] at (2.66,.58)
  {White noise in time\\Endpoint question};

  \draw[rnGray,line width=1pt,dash pattern=on 4pt off 3pt]
  (2.39,.29)--(2.57,.29);
  \node[legend] at (2.66,.29) {Excluded diagonal};
  \draw[threshold] (2.39,.07)--(2.57,.07);
  \node[legend] at (2.66,.07) {White noise lines\\(continued)};
\end{tikzpicture}%
\endgroup%

  \caption[Regular, white and rough noise in one dimension]{Parameter schematic for $d = 1$. The candidate range for
    investigation is $\alpha > 0$, $\alpha_0 > 0$, $\alpha + \alpha_0 < 2$. Blue
    marks the current theorem under its standing covariance assumptions: the
    open regular region, together with the segment white in space, $\alpha = 1$
    and $0 < \alpha_0 < 1$. The hatched rough regions and the dashed segment
    white in time, $\alpha_0 = 1$ and $0 < \alpha < 1$, are questions about
    extending the present results. The axes and the diagonal are excluded,
    including $(1,1)$. The region rough in both time and space does not meet the
    triangle. No solvability claim is made outside the triangle. White in space
    or time means that the corresponding covariance factor is a multiple of a
    Dirac mass. For $d \ge 2$, white noise in space lies on $\alpha = d$,
    outside the open positive candidate triangle.}
  \label{Fig:rough-noise-region}
\end{figure}


Rough spatial covariance is already used in related equations. Balan, Chen and
Ma~\cite{balan.chen.ea:22:parabolic}*{Theorem~1.1} construct the Skorohod
solution of the parabolic Anderson model for their class of initial data,
with $H_0 \in (1/2, 1)$, $0 < H < 1/2$ and
$H_0 + H > 3/4$. In the notation above, this gives $0 < \alpha_0 < 1$, $1 <
\alpha < 2$ and $\alpha_0 + \alpha < 5/2$. Their frequency estimates and bounds
for ratios of Gamma functions suggest methods to investigate. Their estimates
for the heat kernel and orthogonal chaos expansion do not directly treat the
oscillatory wave kernel or the self-pairings in an ordinary Stratonovich
expansion. Related results for wave equations with rough spatial noise are
given by Song, Song and Xu~\cite{song.song.ea:20:fractional}*{Theorem~3.2} and
Liu, Hu and Wang~\cite{liu.hu.ea:22:nonlinear}; their parameter ranges and
definitions of the stochastic integral must also be retained.

The present proof uses nonnegative physical covariance in its pairing bounds
and comparisons of masked moments. It also uses Young integration to construct
the wave evolution. These arguments require further work for the proposed
spectral model. Their limitations do not by themselves imply that the model is
ill posed. At each regularization, noise enters the velocity equation, while
displacement is its time integral. This structure may permit a different
construction. Even after a construction, the sharp moment asymptotics would
remain a separate question. The rough temporal range requires particular
care. The formal continuation of the expression in~\eqref{E:kappa} gives
\[
  \frac{4-\alpha-\alpha_0}{3-\alpha}
  = 1+\frac{1-\alpha_0}{3-\alpha}.
\]
It equals one at $\alpha_0 = 1$ and falls below one when $\alpha_0 > 1$.
For $\alpha_0 \ge 1$, the factor $\delta^{(1-\alpha_0)/2}$ used in the
current block error estimate no longer tends to zero. This factor still
vanishes in the proposed rough spatial range with $\alpha_0 < 1$. We leave
both the appropriate moment scale and its limiting constant open in these
extensions.
